\documentclass{beamer}
\usetheme{Boadilla}
\title{Advanced overlays}
\author{O. Overlay}
\date{October 2026}
\begin{document}
\begin{frame}\titlepage\end{frame}
\begin{frame}{Incremental lists}
  \begin{itemize}[<+->]
    \item One.
    \item Two.
    \item Three.
  \end{itemize}
  \begin{enumerate}
    \item<+-| alert@+> Alert as it appears.
    \item<+-| alert@+> And the next one.
  \end{enumerate}
\end{frame}
\begin{frame}{Pause, only, uncover}
  First line. \pause
  Second line after a pause. \pause

  \only<3>{Only on slide three.}\only<4>{Only on slide four, and this sentence is longer.}

  \uncover<3->{Uncovered from slide three.}

  \visible<2>{Visible on two.} \invisible<2>{Invisible on two.}
\end{frame}
\begin{frame}{Alternatives and temporal}
  \alt<2>{This is the second slide.}{This is not the second slide.}

  \temporal<3>{Before three.}{Exactly three.}{After three.}

  \alert<4>{Alerted on four.} \structure<2-3>{Structure on two and three.}

  \onslide<5>{The end.}
\end{frame}
\begin{frame}{Overlay area and actions}
  \begin{overlayarea}{\textwidth}{3cm}
    \only<1>{A short text.}
    \only<2>{A much longer text that fills more than one line, so that
    the overlay area is what keeps the rest of the slide in place.}
  \end{overlayarea}
  \begin{actionenv}<2->
    Inside an action environment.
  \end{actionenv}
  \action<3->{An action.}
\end{frame}
\begin{frame}[label=transitions]{Transitions}
  \transdissolve<2>
  \transwipe<3>[direction=90]
  Slides with transitions. \pause More. \pause Last.
\end{frame}
\begin{frame}{Buttons and links}
  \hyperlink{transitions}{\beamergotobutton{Back to transitions}}
  \beamerskipbutton{Skip}
  \beamerreturnbutton{Return}
  \hypertarget<2>{target}{A target on slide two.}
  \pause
  \hyperlink{target<2>}{Go to the target}
\end{frame}
\begin{frame}{Labelled overlays}
  \begin{itemize}
    \item<+(1)-> Offset start.
    \item<.-> Same slide as before.
    \item<+-> Next.
  \end{itemize}
  \uncover<4>{Last.}
\end{frame}
\end{document}
